\chapter{医学影像应用}

\paragraph*{本章阅读锚点}
医学场景中的 3DGS 不能只看“新视角是否清晰”。不同模态的成像物理完全不同：内窥镜是窄视野、强反光、器械遮挡和软组织非刚性运动；CT 与 X 射线是从投影恢复体密度；MRI 可能是带物理约束的时空速度场；医学可视化则关心体数据能否在 VR 或自然语言交互中被实时探索。因此本章的读法是先找任务锚点，再看高斯表示被改到哪里：是初始化点云、变形场、抗锯齿滤波、光照嵌入、X 射线 rasterization、物理归一化、还是交互式编辑。

\section{内窥镜场景}

内窥镜路线的基础锚点是 EndoNeRF/EndoSurf 到 EndoGaussian 一类手术场景重建方法：输入是单目或双目手术视频，目标是在组织形变、器械遮挡和镜头照明变化下重建可实时渲染的表面。原始 3DGS 的显式高斯和快速渲染很适合手术导航，但它默认场景静态、光照相对稳定、相机视角较完整；内窥镜方法的改动集中在三点：更好的动态变形、更稳的初始化与遮挡处理，以及针对高频细节和曝光变化的专门约束。



% METHOD: 2405.17872
\paragraph{HFGS（空间与时间高频强调）（2024.05）}
HFGS~\cite{2405_17872} 继承的是 4DGS 内窥镜重建，但它从频率角度解释欠重建：普通像素 L1 或 SSIM 损失在大面积平滑组织上梯度占主导，细小血管、组织边界和快速运动区域的梯度可能被平均掉，导致高斯 densification 和变形优化不足。HFGS 的改变是分别强调空间高频和时间高频，让高斯分裂和运动建模更多关注细节区域。

具体做法中，Spatial High-Frequency reconstruction 先对真实图像和渲染图做傅里叶分析，用高频掩码提取边缘和细纹理分量，再把高频差异作为额外损失或 densification 引导，使高斯在结构细节处更容易被分裂和优化。Temporal High-Frequency reconstruction 使用光流先验识别运动密集区域，对这些区域的变形监督和重建损失加权，让动态组织的快速变化不被静态背景淹没。方法还使用 Depth Anything 等深度先验，并以概率方式建模深度不确定性，避免错误深度过强约束。相对 Endo-4DGX 的光照建模，HFGS 主要解决“细节和运动被平均损失忽视”的问题；相对 SAGS 的滤波抗锯齿，它更强调训练监督中的频率成分。
\textbf{实现细节：} HFGS（2024.05） 从频率角度出发，认为内窥镜动态场景的信息集中分布在特定频段。为此，HFGS 分别引入空间高频强调重建（SHF）和时间高频强调重建（THF），利用光流先验聚焦于运动密集区域，在频域层面选择性增强关键信息。在 Baseline 做法中，现有 4DGS 内窥镜方法（如 EndoGS、EndoGaussian）采用变形场网络建模组织动态，通过多分辨率 HexPlane 编码时空特征。然而在高斯致密化过程中，标准逐像素 $L_1$ 损失的平均梯度在重建不充分区域依然很小。这导致高斯基元无法触发分裂或克隆，产生持续性的欠重建，表现为渲染图像中的模糊和伪影。HFGS 的核心改进是从空间和时间两个频率维度同时解决欠重建问题：空间维度通过 SHF 模块在频域增强静态细节重建；时间维度通过 THF 模块利用光流先验增强动态区域的重建保真度。在实现上，SHF 模块对 ground truth 图像经 FFT（快速傅里叶变换） 获取幅度谱，引入二值掩码提取高频幅度，再经逆 FFT 回到空间域得到高频图像，最终以高频图像为权重对逐像素误差加权，使优化重点聚焦于空间高频（边缘、纹理）区域。THF 模块利用预训练光流预测器（基于 UniMatch）分别对渲染图像和 ground truth 图像预测相邻帧间光流，以 Charbonnier 损失和 Census 损失之和约束两者一致，迫使运动剧烈区域的渲染质量提升。设计动机在于：标准 $L_1$ 损失的均匀惩罚使高频结构和低频平坦区域获得相似权重，但高频区域正是欠重建最明显的位置；变形场网络对快速运动的组织重建不足，而光流天然编码了帧间运动幅度，可以自然定位运动密集区域。两个模块仅在训练阶段施加正则化，不引入额外推理开销。
\textbf{方法信息：} 代码：\url{https://github.com/zhaohaoyu376/hfgs}。
\textbf{设计动机与边界：} 相对原始 3DGS 以独立高斯和图像重建误差为主的设定，HFGS 把问题落在以下可实现的改动上：HFGS包含两个核心模块。（1）空间高频强调重建（SHF）：观察到标准像素级L1损失的平均梯度可能很小，导致高斯致密化不充分（高斯不被分裂）。SHF使用FFT提取图像的空间高频分量，在高频域计算渲染图像与GT的差异，引导高斯致密化关注结构细节。具体：对GT图像做FFT，用二值掩码提取高频振幅，IFFT回空间域计算高频损失。（2）时间高频强调重建（THF）：利用光流先验识别运动密集区域（时间高频分量），在优化中对这些区域施加更高权重。使用预训练光流网络预测GT光流作为监督信号，引导变形场学习动态组织的精确运动。（3）深度概率模型：利用Depth-Anything提供深度先验，用概率模型建模深度不确定性。其设计目的不是单独增加一个模块，而是用上述表示、监督或渲染约束消除原流程在目标任务中的具体失配。从空间和时间频率角度分析并解决3DGS的欠重建问题。SHF通过FFT提取高频分量引导高斯致密化；THF利用光流先验聚焦运动密集区域的优化。在两个广泛使用的基准数据集上达到SOTA渲染质量。适用时还应保留论文给出的边界：依赖光流估计和深度估计的准确性；额外频率计算增加训练开销；仅在手术内窥镜场景验证。
% METHOD: 2503.22437
\paragraph{EndoLRMGS（器械与组织分治重建）（2025.03）}
EndoLRMGS~\cite{2503_22437} 的锚点是 EndoGaussian 式的组织表面重建，但它指出手术场景里有两类对象不应混在同一个表示里：软组织是可变形表面，手术器械则是有闭合形状的刚性或近刚性物体。直接用 3DGS 重建全场景时，器械会遮挡组织，且器械背面和完整形状很难从内窥镜视角恢复。EndoLRMGS 的改变是采用分治策略：组织继续用 Gaussian Splatting 做表面重建，器械则用 Large Reconstruction Model 生成完整 3D 模型，再通过投影优化放回场景。

具体做法先用 DEVA 做零样本分割与跟踪，为每个手术器械得到时序一致掩码。组织分支用掩码引导渲染，把器械遮挡区域从组织监督中剔除，避免把器械颜色和几何吸收到组织高斯里。器械分支把每个工具送入 LRM：图像经 DINOv2 变成 patch token，再投影到三平面表示，MLP 预测 RGB 和密度，最后用 marching cubes 提取闭合网格。由于 LRM 输出的尺度和位置未必与内窥镜场景一致，方法引入 Orthographic-Perspective joint Projection Optimization，同时利用正交投影的稳定轮廓和透视投影的相机几何，优化工具模型的位置、尺度和深度一致性。它相对 EndoGaussian 的核心差异不是换一个渲染损失，而是把“组织表面”和“器械闭合物体”拆成不同三维先验，最后再组合成完整手术场景。
\textbf{实现细节：} 设计动机在于：内窥镜场景的纹理稀疏和结构复杂使传统 SfM 频繁失败，SurGSplat 通过交替优化位姿和场景（而非 SfM 先行、重建后行）实现了端到端的术中重建，特别适合血管结构等关键特征的高保真重建。EndoLRMGS（2025.03） 将大重建模型（LRM）与 GS 结合用于完整手术场景重建，解决 ColonSplat 等纯 GS 方法无法处理器械等封闭物体的 3D 完整性问题。ColonSplat 仅用高斯基元重建组织表面，对于手术器械等需要完整 3D 几何（包括不可见面）的物体无能为力。EndoLRMGS 的核心差异是引入混合表示策略：EndoGaussian 重建组织（表面物体），LRM 重建器械（封闭物体）。具体而言，首先使用 DEVA 分割网络对每帧中的手术器械生成唯一掩码，区分组织与器械；对于组织区域，沿用 EndoGaussian 的反投影初始化和高斯渲染方案，以工具掩码排除器械区域后训练组织的高斯表示；对于每个器械，利用 DINOv2（Distillation with NO labels v2，自监督视觉Transformer） 将分割掩码转换为 patch 特征 token ，投影为三平面表示，再经 MLP 预测 384$^{}3$ 分辨率的 3D 密度场和颜色场，最终通过 Marching Cubes 提取完整水密网格。LRM 生成的器械 3D 模型在固定网格内缺乏物理尺度和空间定位，为此 EndoLRMGS 提出正交透视联合投影优化（OPjPO）：利用 GS 重建的组织点云构建正交投影矩阵，将器械和掩码区域的组织投影到 2D 平面（消除透视畸变），通过面积比估计器械的缩放因子 $\gamma = A_{mask}/A_{tool}$，然后在透视投影下通过最大化器械投影与实际掩码的 IoU 沿 $x$、$y$、$z$ 三轴精调器械位置。设计动机在于：术中场景同时包含变形组织和刚性器械，两者需要截然不同的 3D 表示——高斯基元擅长开放表面渲染，但无法保证封闭物体的拓扑完整性；LRM 擅长从单视图生成完整 3D 模型，但无法直接与 GS 场景配准。EndoLRMGS 通过 OPjPO 弥合两种表示的尺度-位置鸿沟，输出组织-器械一体的完整手术场景 3D 模型。
\textbf{方法信息：} 代码：\url{https://github.com/michaelwanggo/endolrmgs}。



% METHOD: 2506.23308
\paragraph{Endo-4DGX（光照感知动态重建）（2025.06）}
Endo-4DGX~\cite{2506_23308} 的锚点是 4DGS 在可变形手术场景中的应用，但它把问题焦点放在内窥镜最常见的光照失败上：欠曝光、过曝光、器械反光和光源方向变化。普通 4DGS 把外观变化都当作颜色或几何去拟合，容易把曝光错误固化进高斯；Endo-4DGX 的改变是显式引入 illumination embedding，把亮/暗图像和空间区域的校正纳入渲染流程。

具体流程先通过前向-逆向光照估计算法为每帧得到光照先验，并把图像划分成亮、暗或正常曝光状态。每个训练视图有可学习光照嵌入，Region-Aware Enhancement 模块在 Gaussian 级别对局部子区域做亮度增强；为避免亮区和暗区统计差异太大导致训练不平衡，它用两个独立网络分别处理过曝和欠曝校正。Spatial-Aware Adjustment 模块则学习跨视图一致的亮度调整，避免逐帧修复带来时间闪烁。最后，曝光控制损失把不利曝光拉回正常范围，训练目标同时包含重建质量和照明恢复。公开 README 将其定位为 robust endoscopic reconstruction and illumination correction，说明它不是单独的 2D 图像增强器，而是把光照校正嵌入 4DGS 训练和渲染。相对 SAGS 处理动态 aliasing，Endo-4DGX 处理的是动态光照与重建耦合。
\textbf{实现细节：} 由于直接估计指数 $k$ 受相机自动增益干扰且高斯基元几何属性不精确，ColIAGS 不显式估计 $k$，而是通过两项改进隐式建模光照衰减：(1) 余弦嵌入几何建模（IMVE）——在锚基高斯的几何预测 MLP 中，将相机到锚点的方向 $v$ 经余弦基函数 $\Psi(v)$ 投射到高维空间，使几何属性（位置、缩放、旋转）的预测能够感知视角变化，间接提升几何精度以满足光照建模的表面假设；(2) 光照衰减外观建模（IMIA）——在颜色预测 MLP 中，将 $\Psi(v)$ 和距离 $d$ 同时作为输入，由 MLP 隐式学习 $\cos^k\theta$ 衰减行为。这一设计的关键差异在于：ST-Endo4DGS 用柱面谐函数统一编码视角和时间外观，而 ColIAGS 将光照衰减从外观中解耦为独立的隐式分量，并通过物理近似将两个衰减源（径向距离、角度依赖）统一为单一角度参数化。Endo-4DGX（2025.06） 处理极端光照（低光和过曝光）条件下的重建，在 4DGS 基线上引入层次化光照校正管线。与 ColIAGS 从物理模型出发建模稳态光照衰减不同，Endo-4DGX 聚焦于内窥镜手术中动态变化的极端光照场景——术中探头的移动导致同一场景在不同帧中出现低光和过曝光交替出现。
\textbf{方法信息：} 代码：\url{https://github.com/lastbasket/endo-4dgx}。



% METHOD: 2510.27318
\paragraph{SAGS（动态手术场景的自适应抗锯齿）（2025.10）}
SAGS~\cite{2510_27318} 继承的是动态 4DGS 与 Mip-Splatting 的抗锯齿思路。手术内窥镜中组织运动快、纹理细、反光强，普通 4DGS 在高斯投影尺度和变形估计不稳定时会产生闪烁、边界锯齿和时间不一致。SAGS 的改变是把抗锯齿从普通静态视角推广到动态手术场景：一边用注意力驱动的 4D 变形解码器建模组织运动，一边在 3D 高斯空间和 2D 渲染空间同时做滤波。

具体做法包括三个模块。第一，动态加权 4D 变形解码器用多头注意力为不同时间步或运动分量分配权重，让每个高斯随时间的位移更适应局部组织运动，而不是使用固定形式的变形场。第二，3D 平滑滤波器在高斯原语空间抑制过细、过尖或不稳定的几何变化，减少三维层面的 aliasing。第三，2D Mip 滤波器在投影后按屏幕尺度处理高斯 footprint，避免小尺度高斯在像素栅格上产生高频闪烁。它还结合单目或双目深度先验初始化点云，以缓解内窥镜视角稀疏的问题。相对 HFGS 更偏频域高频监督，SAGS 更偏显示和几何层面的抗锯齿处理，目标是在 EndoNeRF 与 SCARED 这类动态手术基准上同时改善 PSNR、SSIM、LPIPS 和视觉稳定性。
\textbf{实现细节：} SAGS（2025.10） 引入注意力驱动的动态加权 4D 变形解码器，从变形场网络架构层面改进 ColonSplat 和 ST-Endo4DGS 使用的固定权重 MLP 解码器。SAGS 的核心差异在于提出自适应加权变形解码器：将多头自注意力机制与 MLP 分支通过可学习的动态权重 $\lambda_1$、$\lambda_2$ 进行融合，使模型在内窥镜场景的大尺度组织运动和细微纹理变化之间自适应切换——自注意力捕获长程依赖保持全局几何一致性，MLP 处理局部非线性变形保留精细细节。解码过程采用双阶段仿射变换（Affine）稳定特征缩放，残差输出通过 $\lambda_1$、$\lambda_2$ 加权后直接更新高斯的位置、旋转、缩放、球谐系数和不透明度。此外，SAGS 针对内窥镜场景特有的视觉伪影（镜面反射高光、窄立体基线和快速非刚性运动引起的锯齿和闪烁）引入抗锯齿处理：3D 平滑滤波器根据 Nyquist 采样定理对每个高斯基元施加低通正则化 $G_k(x)_{reg} = (G_k * G_{low})(x)$ 抑制体积域高频伪影；2D Mip 滤波器在屏幕空间对每个像素区域进行盒式滤波近似，消除缩放视图或相机距离变化时的锯齿。
\textbf{方法信息：} 代码：\url{https://github.com/eververas/sags}。



\section{CT与X射线断层重建}

CT 与 X 射线重建的锚点不是新视角合成，而是断层成像的投影物理：测量值是 X 射线穿过物体后的衰减积分，目标是从稀疏或动态投影恢复三维密度场。标准 3DGS 的颜色、不透明度和 alpha compositing 不能直接代表 X 射线衰减，因此这一节的方法都要改写渲染方程。R$^2$-Gaussian 是本路线的关键锚点，它把高斯从 radiance primitive 改成 radiative density kernel，并推导了 X 射线投影中的积分校正。


% METHOD: 2405.20693
\paragraph{R$^2$-Gaussian（稀疏视角断层的辐射高斯）（2024.05）}
R$^2$-Gaussian~\cite{2405_20693} 是医学断层路线的关键锚点。原始 3DGS 的高斯表示的是可见表面的颜色和不透明度，渲染依赖从近到远的 alpha compositing；CT/X-ray 的像素却来自射线穿过体内组织后的衰减积分。R$^2$-Gaussian 的改变是把高斯从 radiance primitive 改成 radiative Gaussian：每个基元只描述局部体密度，参数包括中心密度、位置和协方差，不再包含视角相关颜色。整个体密度由所有高斯密度核相加得到，投影图像则由 X 射线沿路径积分这些密度产生。

具体管线先用 FDK 等传统 CT 算法从稀疏投影得到一个粗糙体积，再从该体积初始化一组 radiative Gaussians。训练时不能直接套用 3DGS 的相机投影公式，因为把三维高斯压到二维探测器平面时，标准 splatting 会遗漏与协方差相关的积分缩放项；这个 integration bias 会让同一体密度从不同视角看起来不一致。论文因此重新推导 X-ray rasterization：三维密度高斯沿射线积分后变成探测器平面上的二维高斯贡献，并通过 refactored projection 校正缩放因子。优化目标包括投影图像与实测 X-ray 的一致性、密度/几何正则，以及用于输出体数据的可微体素化；代码中也提供 CUDA voxelizer，把连续高斯密度查询为可评估的 3D volume。相对 SAX-NeRF/INR 断层方法，它仍是按病例优化，但不需要每条射线密集采样 MLP，因此能把稀疏视角 CT 重建压到分钟级；后续 IXGS 和 DIGS 基本都是在这个“辐射高斯 + X-ray rasterization”锚点上处理真实术中域差异或动态 CBCT。

\textbf{锚点细化：} R$^2$-Gaussian：是首个将~3DGS~用于稀疏视角层析重建的框架。其核心创新在于提出了"辐射高斯"（Radiative Gaussian）表示——与标准~3DGS~的高斯包含视角相关颜色（SH~系数）不同，辐射高斯移除了视角依赖的颜色分量，因为~X~射线衰减仅取决于各向同性密度。每个辐射高斯基元定义一个局部~Gaussian~形密度场：$G_i^3(x | \rho_i, p_i, \Sigma_i) = \rho_i \exp(-\frac{1}{2}(x - p_i)^T \Sigma_i^{-1}(x - p_i))$，其中~$\rho_i$~为中心密度、$p_i$~为位置、$\Sigma_i = R_i S_i S_i^T R_i^T$~为协方差。整体密度通过求和计算：$\mu(x) = \sum_{i=1}^{M} G_i^3(x)$。
\textbf{实现细节：} R$^{}2$-Gaussian（2024.05） 是首个将 3DGS 用于稀疏视角层析重建的框架。其核心创新在于提出了"辐射高斯"（Radiative Gaussian）表示——与标准 3DGS 的高斯包含视角相关颜色（SH 系数）不同，辐射高斯移除了视角依赖的颜色分量，因为 X 射线衰减仅取决于各向同性密度。每个辐射高斯基元定义一个局部 Gaussian 形密度场：$G_i^{}3(x | \rho_i, p_i, \Sigma_i) = \rho_i \exp(-\frac{1}{2}(x - p_i)^T \Sigma_i^{-1}(x - p_i))$，其中 $\rho_i$ 为中心密度、$p_i$ 为位置、$\Sigma_i = R_i S_i S_i^T R_i^T$ 为协方差。整体密度通过求和计算：$\mu(x) = \sum_{i=1}^{M} G_i^{}3(x)$。在 X 射线光栅化推导中，R$^{}2$-Gaussian 发现了一个关键性质：归一化 3D 高斯分布沿一个坐标轴的积分直接得到归一化 2D 高斯分布。因此，X 射线投影可以通过简单求和 2D 高斯来渲染（而非 alpha-compositing）：$I_r(r) = \sum_{i=1}^{M} G_i^{}2(\tilde{x} | \tilde{\rho}_i, \tilde{p}_i, \tilde{\Sigma}_i)$，其中 $\tilde{\rho}_i = \rho_i \cdot (2\pi)^{3/2}|\Sigma_i|^{1/2}$ 是投影后的 2D 高斯中心密度。更关键的是，R$^{}2$-Gaussian 发现了标准 3DGS 中未知的"积分偏差"（integration bias）——3DGS 实际上学习的是 2D 图像平面上的积分密度而非 3D 空间中的实际密度，这导致不同视角恢复出不同的 3D 密度值，违反了 X 射线衰减的各向同性本质。R$^{}2$-Gaussian 通过将实际 3D 密度分配给高斯基元并正向计算 2D 投影，从根本上解决了这一问题。此外，该方法开发了可微体素化器（voxelizer）用于 3D 密度体积查询和全变分（TV）正则化，初始化采用 FDK 解析重建的快速低质量体积，在 4 分钟内完成高质量重建，比 NeRF 方法快 12$\times$。
\textbf{方法信息：} 代码已公开：\url{https://github.com/Ruyi-Zha/r2_gaussian}，仓库约 291 stars；Semantic Scholar 约 73 次引用。



% METHOD: 2504.14699
\paragraph{IXGS（稀疏任意姿态真实 X 射线术中重建）（2025.04）}
IXGS~\cite{2504_14699} 也是从 R$^2$-Gaussian 出发，但它面对的是更贴近手术导航的真实 X-ray：视角稀疏、姿态任意、真实 X-ray 与合成 DRR 存在明显域差异。标准 R$^2$-Gaussian 在理想几何和较干净投影上表现好，但术中 C-arm 或透视图像的对比度、散射、遮挡和采集轨迹都更复杂。IXGS 的改变是在辐射高斯重建前加入 anatomy-guided radiographic standardization，把真实 X-ray 先转换到更接近合成 DRR 的标准域。

具体管线先采集配对的 CT、真实 X-ray 和由 CT 生成的 DRR，并完成几何标定。然后用 Pix2Pix 式图像到图像模型学习从真实 X-ray 到标准化 DRR 域的转换，使骨结构更清晰、辐射风格更一致。接着进行裁剪、体积归一化和相机几何预处理，再用 IXGS 框架执行稀疏投影三维重建。它不依赖大规模跨病人预训练，而是按实例学习来适应新患者解剖。相对 DIGS，它不处理时间运动；相对 R$^2$-Gaussian，它的增量在真实术中成像域适配和任意姿态稀疏视图。专家外科评估显示，约 20 到 30 个标准化视图即可达到可用于导航判断的重建质量，这说明 3DGS 断层方法从合成 benchmark 走向真实手术场景时，域标准化和几何预处理与核心 splatting 同样重要。
\textbf{实现细节：} IXGS（2025.04） 将 3DGS 扩展到术中稀疏任意姿态 X 射线的 3D 重建。\textit{与 R$^{2}$-Gaussian 的差异。}R$^{2}$-Gaussian 假设 X 射线投影的几何参数（视角、姿态）是精确已知的。术中场景（如脊柱手术中的 C 臂 X 射线）面临两个额外挑战：投影姿态可能是任意且不规则的，不同视角间的视觉一致性差（不同曝光、不同组织重叠）。IXGS 在 R$^{2}$-Gaussian 框架上引入解剖引导的射线标准 化（anatomy-guided radiographic standardization）：使用风格迁移统一不同视角间的视觉外观，提升跨视角一致性从而增强重建质量。无需预训练，天生适应新患者和解剖结构。
\textbf{实现补充：} 实现上还要核对输入、输出、参数、特征、初始化、训练损失、正则、渲染或推理接口，以及优化中的剪枝、分裂、采样、编码或解码步骤；这些抓手决定该方法到底改了 3DGS 管线的哪一环。
\textbf{方法信息：} 代码：\url{https://github.com/mrmonk3y/ixgs}。
\textbf{设计动机与边界：} 相对原始 3DGS 以独立高斯和图像重建误差为主的设定，IXGS-Intraoperative 3D Reconstruction from Sparse, Arbitrarily Posed Real X-rays 把问题落在以下可实现的改动上：IXGS扩展了R2-Gaussian溅射框架以处理术中X射线数据。核心挑战在于真实X射线与合成DRR之间存在域差异。提出解剖引导放射标准化：使用Pix2Pix模型训练于配对的真实X射线和合成DRR数据，将真实X射线转换到标准化域。管道包括：(1)配对的CT和X射线数据采集与标定；(2)Pix2Pix风格的域转换；(3)几何预处理（图像裁剪和体积归一化）；(4)基于IXGS框架的体3D重建。框架无需预训练，利用实例学习方法避免大量标注需求，天然适应新患者解剖结构。其设计目的不是单独增加一个模块，而是用上述表示、监督或渲染约束消除原流程在目标任务中的具体失配。首次将高斯溅射应用于术中稀疏任意姿态真实X射线的3D重建，引入基于风格变换的解剖引导放射标准化步骤提高视觉一致性，无需预训练即可适应新患者解剖结构，经专家外科评估确认临床实用性。适用时还应保留论文给出的边界：重建精度不如CBCT；需要至少20-30个视图才能获得可接受的重建质量；依赖配对训练数据用于标准化。
% METHOD: 2506.22280
\paragraph{DIGS（动态 CBCT 的变形感知 4DGS）（2025.06）}
DIGS~\cite{2506_22280} 继承 R$^2$-Gaussian 的 CT 投影建模，但把静态体密度扩展到动态 CBCT。临床 4D CBCT 常要处理呼吸运动，传统做法按呼吸相位分箱后分别重建，容易增加剂量并受相位排序错误影响；DIGS 的改变是直接在逐投影时间点上建模动态体积，用低秩自由形变场驱动高斯随时间演化。

具体做法先用一组基准高斯表示参考时刻的衰减场，再引入 Free-Form Deformation 空间基函数描述呼吸运动。不同于 HexPlane 或通用 4DGS 中分别预测高斯位置、尺度和旋转的变化，DIGS 提出 deformation-informed 框架：先估计统一的 deformation vector field，再用该向量场一致更新高斯均值，并通过局部线性化把协方差的变化也从同一变形中推导出来。这样，高斯的几何属性不会彼此独立漂移，重建出的运动场也更有物理意义。它还利用低秩运动模型减少参数量，支持每个投影独立对应一个时间状态，而不强依赖规则呼吸相位。相对 R$^2$-Gaussian，DIGS 的主要增量是“静态密度核加时间变形”；相对普通 4DGS，它的渲染仍遵循 X 射线投影物理，而不是 RGB alpha compositing。
\textbf{设计动机与边界：} 相对原始 3DGS 以独立高斯和图像重建误差为主的设定，DIGS 把问题落在以下可实现的改动上：核心创新：(1)FFD空间基函数：提出基于自由变形的显式运动表示替代HexPlane等隐式表示，计算效率更高。FFD通过控制网格的变形来建模器官的呼吸运动。(2)变形感知(DI)框架：不同于现有方法独立修改高斯核的均值/缩放/旋转，DI框架通过统一的变形向量场(DVF)一致更新所有高斯属性。均值位置直接通过DVF偏移，协方差矩阵通过局部线性近似DVF推导。保证高斯运动的空间一致性，可推导与渲染图像一致的连续DVF。(3)低秩运动模型：结合FFD的低秩特性，减少运动参数数量。(4)逐投影重建：不需相位排序，每个投影独立重建，捕捉不规则呼吸运动。为了消除原始 3DGS 在该目标任务中的具体失配，方法把上述表示、监督或渲染约束作为可学习和可执行的改动，而不是单独堆叠模块。提出基于自由变形(FFD)的空间基函数和变形感知框架，统一约束高斯均值/缩放/旋转的时序演化，实现高效动态CBCT重建，图像质量优于HexPlane且速度快6倍。适用时还应保留论文给出的边界：FFD控制网格分辨率影响精度；对大幅非刚性变形的建模能力有限；依赖呼吸信号的准确性。
\textbf{端到端流程：} DIGS 的输入是逐投影医学观测及呼吸相关的形变信息。首先在规范空间初始化高斯和 FFD 控制网格；FFD 产生的连续 DVF 同时更新均值、尺度和旋转，其中协方差由 DVF 的局部线性近似传播，而不是分别预测各属性。训练按每个投影独立重建，无需相位排序，以投影重建误差驱动高斯和低秩 FFD 参数；推理时对给定时间/呼吸状态评估 DVF，变形高斯后通过标准 splatting 渲染对应投影或体积外观。
\textbf{方法信息：} 代码：\url{https://github.com/yuliang-huang/digs}。



\section{手术场景重建}

这一节与内窥镜动态重建有交叉，但更强调完整手术系统：如何用深度先验、器械 mask、硬件后端和边缘部署把 3DGS 推向手术室实时管线。这里的锚点是“临床可用的速度和稳定性”，不只是研究基准上的平均 PSNR。


% METHOD: 2405.17835
\paragraph{Deform3DGS（快速手术非刚性重建）（2024.05）}
Deform3DGS~\cite{2405_17835} 是较早把 3DGS 系统性引入手术动态场景的工作之一。它针对 EndoNeRF 类方法训练慢、渲染慢，以及普通 4DGS 变形解码器开销大的问题，做了两处关键改变：用多视角深度融合替代稀疏 SfM 初始化，用显式基函数替代特征平面加 MLP 的时间变形表示。

具体做法的第一部分是 Multi-view Aware Point Fusion。它利用立体深度或深度先验和视觉里程计，把多帧深度图融合成更密集、更贴近组织表面的初始点云，再初始化高斯。这样训练一开始就有较好的几何覆盖，不必完全依赖后续 densification。第二部分是 Flexible Deformation Modeling。对于每个高斯，方法不用 HexPlane 查询时空特征再经 MLP 解码，而是用傅里叶基、多项式基等时间基函数的线性组合表示运动轨迹，并学习每个高斯的组合权重；必要时再加入可学习规范化位移处理超出基函数表达的非刚性形变。训练仍通过图像重建和深度监督端到端优化。它相对 HFGS/SAGS 更早、更轻量，重点是用合适初始化和低成本变形把手术场景训练压到分钟级，并保持实时渲染。
\textbf{实现细节：} Deform3DGS（2024.05） 首次将 3DGS 引入手术场景的快速重建。它修改了 3DGS 的变形建模环节，用灵活变形建模（FDM）替代 HexPlane 隐式变形场和傅里叶/多项式显式变形；3D 高斯基元的光照/透明度属性和可微渲染保持不变。Baseline 做法及其局限。动态场景重建的主流变形建模有两条路线：(1) 基于 HexPlane 的隐式变形场（如 4DGS、LerPlane），将时空信息编码为共享特征平面并通过 MLP 解码时间变形。虽然表示能力强，但特征平面插值和 MLP 推理的计算开销大，难以满足手术场景的实时需求。(2) 基于傅里叶/多项式基函数的显式变形，为每个高斯基元赋予可学习权重 $\omega$ 对固定基函数进行线性组合来回归运动曲线 $\psi(t;\omega)=\sum_{j=1}^{B}\omega_j b_j(t)$。虽然计算高效，但存在两个根本局限：(i) 基函数在给定时间点取值固定、不随训练调整，时间变形被约束为规范运动轨迹的一致表示，模型被迫放弃某些特定运动以保持整体轨迹的连贯性，无法精确捕捉器械-组织交互等精细局部变形；(ii) 某一时刻学到的变形会全局影响整条轨迹，不同查询时间的局部变形表示相互干扰。改了什么：FDM 新机制。FDM 的核心改进是将固定基函数替换为可学习参数化的 Gaussian 基函数——每个基函数的中心 $\mu$ 和宽度 $\sigma$ 均可在训练中优化，使变形曲线不再被固定函数形状束缚，而是自适应地调整每个基函数在时间轴上的"作用区域"和"影响范围"。直觉上，每个 Gaussian 基函数就像时间轴上的一个"可移动、可伸缩的灯泡"，灯泡中心 $\mu$ 决定它照亮哪个时间段，灯泡宽度 $\sigma$ 决定它照亮多宽的时间窗口，而权重 $w$ 决定亮度——多个灯泡叠加即可拟合任意复杂的变形曲线。实现原理。具体地，FDM 为每个高斯基元赋予一组 Gaussian 基函数（论文实验中取 $B=17$），以可学习中心 $\mu_j$ 和方差 $\sigma_j$ 参数化每个基函数： 即每个基函数在时间轴上以 $\mu_j$ 为中心、$\sigma_j$ 为宽度呈钟形分布。
\textbf{方法信息：} 引用很多，Semantic Scholar 约 54 次。
\textbf{设计动机与边界：} 相对原始 3DGS 以独立高斯和图像重建误差为主的设定，Deform3DGS 把问题落在以下可实现的改动上：Deform3DGS包含两个核心创新。（1）多视角感知点云融合（MAPF）：利用立体深度先验和视觉里程计构建初始高斯点云，替代标准SfM稀疏点。通过多视角深度图融合和一致性检查，生成更密集准确的初始点云，显著加速几何重建收敛。（2）灵活变形建模（FDM）：不使用HexPlane等计算密集的特征平面+MLP方案，而是用傅里叶和多项式基函数的线性组合来学习每个高斯的时间运动曲线。每个高斯有一组可学习权重，线性组合基函数生成运动轨迹：δ(t;θ)=Σw\_j*b\_j(t)。FDM还引入可学习的规范化位移机制，处理非刚性变形超出基函数表达范围的情况。整个框架通过深度监督和图像重建损失端到端优化，无需复杂的特征平面计算。其设计目的不是单独增加一个模块，而是用上述表示、监督或渲染约束消除原流程在目标任务中的具体失配。提出Deform3DGS，首次在手术场景中实现3DGS的实时渲染和快速训练。MAPF方案利用深度先验改善点云初始化，FDM方案用基函数学习每个高斯的时间变形曲线，无需特征平面插值和MLP解码。实现PSNR 37.90、渲染速度338.8 FPS、训练时间仅1分钟/场景。适用时还应保留论文给出的边界：假设手术场景无器械遮挡；基函数数量限制变形表达能力；长序列累积漂移问题。
% METHOD: 2512.02482
\paragraph{G-SHARP（手术硬件加速实时管线）（2025.12）}
G-SHARP~\cite{2512_02482} 的锚点不是单个新损失，而是可部署手术重建管线。它基于商业兼容的 gsplat 后端和 NVIDIA Holoscan/IGX 边缘硬件，把 3DGS 手术重建从离线研究代码推进到手术室实时可视化。它面对的三个问题是窄视野、器械遮挡和动态组织变形，因此方法同时改初始化、监督 mask 和部署后端。

具体做法采用两阶段 Gaussian Splatting 框架。初始化时，它不是从单帧或稀疏 SfM 开始，而是累积多帧深度和颜色生成更完整点云；在组织模式下，器械 mask 会被用于排除器械区域，从而得到更多干净组织点。训练时，组织模式把器械 mask 同时作用于 GT 和渲染图，损失只从组织区域反传；对所有帧中曾被器械遮挡的“不可见区域”，方法加入 TV 正则，鼓励组织表面做平滑插值。动态部分通过变形网络在精训练阶段作用于组织高斯，并用逐高斯变形表跟踪哪些高斯应变形，避免器械区域干扰组织运动。相对 Deform3DGS，G-SHARP 更强调工程系统：速度、许可证兼容、边缘硬件部署和手术室实时反馈。
\textbf{实现细节：} 由于位姿估计存在噪声，直接用噪声位姿优化纹理会导致模糊伪影。RTL 采用交替优化策略——每轮先精修 N 步位姿参数，再固定位姿优化 M 步纹理参数——逐步消除位姿噪声对纹理质量的影响。G-SHARP（2025.12） 构建首个商业兼容的实时手术重建框架，标志着 3DGS 手术重建从学术概念验证向手术室实际部署的关键一步。G-SHARP 是首个原生构建于 GSplat（Apache-2.0 许可证）可微高斯光栅化器的手术重建管线——GSplat 是 Nerfstudio 项目维护的 PyTorch 兼容后端，提供融合前向/反向传播的高效 CUDA 内核和数值稳定的梯度计算。基于此后端，G-SHARP 实现三重创新：（1）面向 EndoNeRF-Pulling 的变形与掩码方案；（2）两阶段训练管线——粗阶段（ 200 次迭代，约 30 秒）仅初始化几何，精阶段（ 1500 次迭代，约 1.5 分钟）引入时间变形；（3）基于 Holoscan SDK 的实时推理部署。变形网络由 HexPlane 时空特征网格和 MLP 解码器组成——HexPlane 将 4D 时空体分解为 6 个 2D 特征平面（3 个空间平面 XY/XZ/YZ + 3 个时空平面 XT/YT/ZT），空间分辨率 $64^{}3$、时间维度 100，采用多分辨率网格（尺度 $[1,2,4,8]$）捕获不同粒度特征；MLP 解码器（8 层，256 维隐藏单元）从双线性插值特征中预测每个高斯基元在时刻 $t$ 的变形偏移 $\Delta\bm{\mu}$（位置）、$\Delta\bm{s}$（缩放）、$\Delta\bm{q}$（旋转四元数）和 $\Delta\alpha$（不透明度）。训练采用掩码监督策略——工具掩码应用于 GT 和渲染图像后再计算损失，确保梯度仅通过组织区域流动；对所有帧工具掩码取并集生成"不可见掩码"，对始终被遮挡区域施加 TV 正则化鼓励平滑插值。
\textbf{方法信息：} NVIDIA 参与了这篇论文；相关工程资料：\url{https://github.com/nvidia-holoscan/holohub}，当前未确认论文完整代码。
\section{解剖可视化与医学教育}

医学可视化中的 3DGS 主要不是替代 CT/MRI 重建算法，而是把体数据转成可交互、可编辑、可在轻量设备上实时渲染的表示。这里的锚点是 volume rendering 和 cinematic rendering；3DGS 的增量在于可编辑高斯、快速 splatting、自然语言控制和移动 VR 压缩。


% METHOD: 2507.12621
\paragraph{NLI4VolVis（自然语言交互体可视化）（2025.07）}
NLI4VolVis~\cite{2507_12621} 把 3DGS 放在医学和科学体可视化的交互层。传统体可视化需要用户调 transfer function、切面、透明度和颜色映射；NLI4VolVis 的改变是让用户用自然语言查询和编辑体场景，而底层对象用可编辑 3D Gaussians 表示。

具体管线先从体数据渲染多视图图像，再用 LangSplat、SAGD 或 SAM2 等方法做多视图语义分割，把不同器官或结构分成可独立编辑的对象。每个对象训练为 iVR-GS 风格的可编辑高斯组件，支持颜色、不透明度、光照和组合操作。为了让开放词汇查询更稳定，系统用熵引导视角选择挑出信息量高的视图，再用 CLIP/VLM 建立语义理解。交互层由多 Agent LLM 架构处理：一个 agent 解析用户意图，另一个选择工具函数，另一个检查执行结果，并在 VPA 循环中迭代优化场景。相对普通 LangSplat 场景理解，它不只做语义定位，而是把分割、高斯编辑和自然语言控制串成医学体数据探索系统。
\textbf{实现细节：} NLI4VolVis（2025.07） 将体积可视化的交互范式从空间裁剪推进到语义级自然语言交互。\textit{锚点回顾。}ClipGS 的交互模式是几何操作——用户通过移动裁剪平面逐层"解剖"体积数据，系统响应的是空间位置关系。然而，对于非专业用户（如医学生、安全检查人员），他们更希望用自然语言表达意图（如"显示消化系统相关器官"或"把背包里的方形物体变红"），而非手动调整裁剪面或传递函数参数。NLI4VolVis 正是针对这一需求，构建了一个完整的自然语言驱动的体积可视化系统。\textit{系统架构与核心设计。}NLI4VolVis 的管线包含四个阶段。（1）多视角语义分割：从体积数据渲染多视角图像后，采用灵活的分割策略——支持预标注数据集、传递函数分割、体素分割、跨帧对象跟踪（SAM 2）以及直接在 3DGS 上分割（LangSplat、SAGD）等五种方式，每个分割出的语义组件独立训练一个可编辑 3DGS 模型（基于 iVR-GS）。（2）3D 开放词汇场景理解：对每个分割组件，用基于信息熵的视角选择机制从均匀采样的多视角图像中挑选信息量最大的 $k$ 帧（按 alpha 通道熵值 $H(I)=-\sum_{i=1}^{N} p_i \log p_i$ 排序），再通过 CLIP（Contrastive Language-Image Pre-training，对比语言-图像预训练） 模型提取视觉嵌入并取均值 $E_{\mathrm{vision}} = \frac{1}{k}\sum_{j=1}^{k} E_{\mathrm{frame}_j}$，若存在文本标签则融合 $E_{\mathrm{object}} = (E_{\mathrm{vision}} + E_{\mathrm{text}})/2$。用户查询时计算 CLIP 余弦相似度匹配目标对象。
\textbf{实现补充：} 实现上还要核对输入、输出、参数、特征、初始化、训练损失、正则、渲染或推理接口，以及优化中的剪枝、分裂、采样、编码或解码步骤；这些抓手决定该方法到底改了 3DGS 管线的哪一环。
\textbf{方法信息：} 代码：\url{https://github.com/kuangshiai/nli4volvis}。
\textbf{设计动机与边界：} 相对原始 3DGS 以独立高斯和图像重建误差为主的设定，NLI4VolVis 把问题落在以下可实现的改动上：NLI4VolVis的流水线包括：(1)多视图语义分割：使用LangSplat、SAGD或SAM2对体积场景的多视图图像进行语义分割，每个分割对象训练独立的iVR-GS可编辑模型；(2)开放词汇场景理解：通过熵引导视角选择机制，用CLIP模型嵌入信息量最大的帧，实现鲁棒的开放词汇理解；(3)可编辑高斯训练：基于iVR-GS训练每个分割组件，支持颜色、不透明度、光照的实时编辑，以及通过高斯参数拼接的场景组合；(4)分层多Agent系统：LLM Agent利用外部函数调用工具解释用户意图并修改3D场景，遵循VPA循环实现迭代优化。系统还支持语音输入/输出、2D风格化（通过IP2P）和透明Agent操作日志。其设计目的不是单独增加一个模块，而是用上述表示、监督或渲染约束消除原流程在目标任务中的具体失配。提出了NLI4VolVis系统，首次实现了通过自然语言对体积可视化场景进行探索、查询和编辑。系统集成了多视图语义分割、VLM和多Agent LLM架构，基于可编辑3D高斯实现开放词汇对象查询、实时场景编辑、最佳视角选择和2D风格化。通过案例研究和用户研究验证了系统在体积数据探索中的可访问性和可用性提升。适用时还应保留论文给出的边界：依赖外部VLM和LLM的推理能力，系统整体延迟可能较高；2D风格化缺乏多视角一致性；不支持直接在原始体积数据上进行分割；不同分割方法在不同数据集上各有优劣，无通用最优方案。
% METHOD: 2601.19310
\paragraph{ClipGS-VR（移动 VR 体数据切片可视化）（2026.01）}
ClipGS-VR~\cite{2601_19310} 的锚点是 ClipGS 式的体数据 cinematic visualization。原始方法可以产生高质量医学体渲染，但模型状态大、推理重、交互切片受限，不适合消费级移动 VR。ClipGS-VR 的改变是把神经推理结果预烘焙进统一高斯结构，并为任意角度切片设计平滑不透明度调制。

具体做法先在约 200 个均匀采样切片状态下运行 ClipGS 的 Learnable Truncation 和 Adaptive Adjustment Model，把每个状态的高保真输出写入高斯属性。为了压缩，它把高阶球谐系数拆成跨状态共享的全局引用数组，只保留切片平面附近的高保真层基元，最终把单案例数据从 GB 级降到几十 MB。交互切片时，它不再用硬截断直接删除平面一侧高斯，而是根据高斯到切片平面的距离做基于梯度的不透明度调制，使任意角度切片时透明度连续变化，减少锯齿和跳变。相对 NLI4VolVis 的自然语言编辑，ClipGS-VR 更关注“在轻量 VR 设备上实时、平滑地看医学体数据”。
\textbf{实现细节：} 进一步，可学习截断本质上是离散的可见性控制，在裁剪面连续移动时会导致渲染不一致。为此，ClipGS 设计了自适应变形调整模型（Adaptive Deformation Adjustment）：用位置编码+轻量 MLP 将高斯位置和裁剪面参数编码为特征 $f$，再通过三个独立的 MLP 头分别预测位置偏移 $\Delta\mu$、旋转偏移 $\Delta q$ 和缩放偏移 $\Delta s$，将最终渲染改写为仅对可见高斯的加权合成：$C = \sum_{i=1}^{N} M_i T_i \alpha_i G_i'(r; \mu_i + \Delta\mu_i, q_i + \Delta q_i, s_i + \Delta s_i) c_i$。这使得裁剪面移动时高斯能连续变形，消除离散跳变。ClipGS-VR（2026.01） 面向临床场景中外科医生佩戴 VR 头显自由探索 CT 解剖结构的需求，将 ClipGS 从桌面工作站移植到移动 VR 平台（Meta Quest 3）。核心挑战在于：ClipGS 的可学习截断与自适应变形依赖运行时神经网络推理，其 VRAM 占用超过 6\,GB，无法在独立 VR 头显上实时运行。ClipGS-VR 的解决方案分两步。首先，将神经推理预计算烘焙为显式高斯属性——对每个病例均匀采样 200 个切片平面状态，离线计算 LT 和 AAM 的输出并直接嵌入高斯基元，消除运行时推理开销。在压缩策略上，分离球谐系数为全局共享数组（因为所有切片状态共享相同的高阶 SH 系数），并仅保留切片平面处的差异基元而非复制整个体积，将每个病例从约 6\,GB 压缩至平均 40\,MB。其次，将 ClipGS 原有的单轴裁剪平面扩展为6-DoF 任意角度切片：对每个高斯基元，根据其中心到切片平面的符号距离与沿法线方向的投影半径之比计算透明度系数 $\sigma = \mathrm{clamp}(1/2 + (\mathbf{\mu} \cdot \mathbf{n} - c)/(2 \cdot s_{\mathbf{n}}), 0, 1)$，再调制不透明度 $\alpha' = \alpha \cdot \sigma$，使与切片平面相交的高斯按几何穿透深度平滑淡出，消除硬截断的锯齿伪影。
\textbf{技术实现补充：} 进一步，可学习截断本质上是离散的可见性控制，在裁剪面连续移动时会导致渲染不一致。为此，ClipGS 设计了自适应变形调整模型（Adaptive Deformation Adjustment）：用位置编码+轻量 MLP 将高斯位置和裁剪面参数编码为特征 $f$，再通过三个独立的 MLP 头分别预测位置偏移 $\Delta\mu$、旋转偏移 $\Delta q$ 和缩放偏移 $\Delta s$，将最终渲染改写为仅对可见高斯的加权合成：$C = \sum_{i=1}^{N} M_i T_i \alpha_i G_i'(r; \mu_i + \Delta\mu_i, q_i + \Delta q_i, s_i + \Delta s_i) c_i$。这使得裁剪面移动时高斯能连续变形，消除离散跳变。ClipGS-VR（2026.01） 面向临床场景中外科医生佩戴 VR 头显自由探索 CT 解剖结构的需求，将 ClipGS 从桌面工作站移植到移动 VR 平台（Meta Quest 3）。核心挑战在于：ClipGS 的可学习截断与自适应变形依赖运行时神经网络推理，其 VRAM 占用超过 6\,GB，无法在独立 VR 头显上实时运行。ClipGS-VR 的解决方案分两步。首先，将神经推理预计算烘焙为显式高斯属性——对每个病例均匀采样 200 个切片平面状态，离线计算 LT 和 AAM 的输出并直接嵌入高斯基元，消除运行时推理开销。在压缩策略上，分离球谐系数为全局共享数组（因为所有切片状态共享相同的高阶 SH 系数），并仅保留切片平面处的差异基元而非复制整个体积，将每个病例从约 6\,GB 压缩至平均 40\,MB。其次，将 ClipGS 原有的单轴裁剪平面扩展为6-DoF 任意角度切片：对每个高斯基元，根据其中心到切片平面的符号距离与沿法线方向的投影半径之比计算透明度系数 $\sigma = \mathrm{clamp}(1/2 + (\mathbf{\mu} \cdot \mathbf{n} - c)/(2 \cdot s_{\mathbf{n}}), 0, 1)$，再调制不透明度 $\alpha' = \alpha \cdot \sigma$，使与切片平面相交的高斯按几何穿透深度平滑淡出，消除硬截断的锯齿伪影。
\textbf{方法信息：} Microsoft 参与了这篇论文。
\textbf{设计动机与边界：} 相对原始 3DGS 以独立高斯和图像重建误差为主的设定，ClipGS-VR 把问题落在以下可实现的改动上：两项核心创新：(1)数据预处理与压缩：将ClipGS的Learnable Truncation和Adaptive Adjustment Model在200个均匀采样切片状态下预计算，输出直接烘焙到高斯属性中。两项压缩策略：将高阶球谐系数解耦为跨状态共享的全局引用数组；仅保留切片平面处的独特基元（最高保真度层），最终从6GB压缩至40MB。(2)交互式任意角度切片：扩展原始仅支持单轴切割的ClipGS。使用基于梯度的不透明度调制替代硬截断，根据高斯基元到切片平面的距离平滑调制不透明度，避免锯齿伪影和视觉不连续。其设计目的不是单独增加一个模块，而是用上述表示、监督或渲染约束消除原流程在目标任务中的具体失配。提出ClipGS-VR将ClipGS扩展到移动VR平台，支持任意角度切片。(1)将神经推理重构为合并数据集，将200个预计算切片状态的高保真层整合为统一渲染结构，每个案例从6GB压缩至40MB；(2)引入基于梯度的不透明度调制实现任意角度平滑切片渲染。适用时还应保留论文给出的边界：预计算200个切片状态增加了离线处理时间；压缩可能损失极端角度的细节保真度；当前仅验证医学体数据。

\section{MRI重建}

MRI 超分辨率中的 3DGS 不是图像渲染器，而是连续物理场近似器。4D flow MRI 的输入是时空速度场，目标是从低分辨率观测恢复高分辨率血流，同时满足物理约束。这里必须把原始 3DGS 的投影 alpha blending 去掉，改成在时空坐标上直接查询高斯基函数。
% METHOD: 2511.11048
\paragraph{PINGS-X（物理信息归一化高斯超分辨率）（2025.11）}
PINGS-X~\cite{2511_11048} 的锚点是 PINN 和 Physics-Informed Gaussian Splatting。PINN 可以用 Navier-Stokes 等物理方程约束流场，但 MLP 查询慢；普通未归一化高斯求和虽然快，却在远离高斯中心时预测塌缩到零，训练容易靠过大高斯或不断增殖来补偿。PINGS-X 的改变是提出 Normalized Gaussian Splatting：预测值不再是未归一化加权和，而是高斯权重归一化后的凸组合，因此输出有界，并给出收敛保证。

具体做法把 4D 时空坐标映射到速度和压力等物理量，每个高斯带中心、尺度和物理属性。查询某个点时，先计算所有相关高斯的权重，再归一化为凸组合，得到速度/压力预测；这样远离单个中心时不会无意义衰减为零。为适应 4D 时空和多物理量训练，它限制高斯为 axes-aligned，降低旋转参数开销，同时保持理论性质。训练中加入物理信息损失，例如不可压缩流和 Navier-Stokes 相关约束，并用 Gaussian merging 合并相似或退化的高斯，防止多个高斯聚在同一位置造成效率下降。相对图形学 3DGS，PINGS-X 的变化很彻底：没有屏幕投影和 alpha compositing，而是用归一化高斯作为科学计算中的连续函数近似器。
\textbf{实现细节：} PINGS-X（2025.11） 针对 4D Flow MRI 超分辨率提出轴对齐时空高斯表示，引入归一化 3DGS 和收敛保证。\textit{与 3DGSMR 的差异。}3DGSMR 重建 MR 图像的标量强度值（解剖结构），PINGS-X 重建 4D Flow MRI 的速度场向量（血液流速的 3D 分量 $v_x, v_y, v_z$ 随时间 $t$ 变化）——这是一个 4D 向量场重建问题。PINGS-X 引入三个关键创新：(1) 归一化高斯 Splatting——对高斯基元的贡献进行归一化，提供形式化的收敛保证，确保优化过程的稳定性；(2) 轴对齐高斯——将 4D 时空高斯简化为沿各轴独立的 1D 高斯乘积，大幅简化高维数据的训练复杂度同时保持精度和收敛保证；(3) 高斯合并过程——防止退化解并提升计算效率。
\textbf{方法信息：} 代码：\url{https://github.com/spatialailab/pings-x}。
\section{其他医学应用}

这一节收纳跨模态医学影像重建。它们不严格属于内窥镜、X-ray 或 MRI 单一模态，但共同目标是用显式高斯承载可解释、可编辑的医学几何。与自然图像 3DGS 不同，医学场景常面对稀疏切片、多模态灰度、分割边界和解剖连续性；因此这里更关注高斯如何同时服务图像插值、结构重建和临床可视化，而不是单纯提高照片级新视角质量。


% METHOD: 2509.16806
\paragraph{MedGS（多模态 3D 医学影像高斯重建）（2025.09）}
MedGS~\cite{2509_16806} 的锚点是从切片序列恢复三维解剖结构的传统医学影像管线，例如 Marching Cubes、隐式神经表示或分割后网格重建。它的改变是把体数据表示为 3D 空间中的 2D 横截面序列，并用多任务高斯架构同时做灰度帧插值和分割/表面提取，让密集图像合成信号反过来正则化几何。

具体做法中，高斯原语共享一套几何参数，包括位置和可能的时间变形；外观上则分成两个分支，一个用于灰度或多模态图像回归，另一个用于 mask 或分割相关的表面提取。这样，稀疏切片之间的图像插值不仅生成更密集的观测，也帮助几何高斯学习连续、平滑的解剖结构；分割分支再把共享几何转成可用于网格或表面重建的解剖对象。它相对单纯 Marching Cubes 的差异在于，几何不是只从阈值体素中提取，而是在图像合成和分割两个任务共同监督下优化；相对纯 INR，它保留了高斯显式结构，便于编辑和交互。
\textbf{实现细节：} MedGS（2025.09） 提出多模态 3D 医学成像的高斯 Splatting 框架，支持 CT 和 MRI 等多种模态的联合 3D 重建——核心挑战在于不同模态的成像物理差异（CT 测量 X 射线衰减，MRI 测量磁共振信号），需要统一的 3DGS 表示来桥接。
\textbf{实现补充：} 实现上还要核对输入、输出、参数、特征、初始化、训练损失、正则、渲染或推理接口，以及优化中的剪枝、分裂、采样、编码或解码步骤；这些抓手决定该方法到底改了 3DGS 管线的哪一环。
\textbf{方法信息：} 代码：\url{https://github.com/gmum/medgs}。
\textbf{设计动机与边界：} 相对原始 3DGS 以独立高斯和图像重建误差为主的设定，MedGS 把问题落在以下可实现的改动上：MedGS基于VeGaS和Folded-Gaussian分布构建，将体积数据表示为3D空间中的2D横截面序列。核心设计是多任务架构：高斯原语共享统一几何参数（位置、时间变形），同时维护两个独立的颜色分支——一个用于灰度帧插值回归，另一个用于基于掩码的表面提取。这种共享几何+双外观设计使得密集的图像合成信号能够正则化几何学习，即使从稀疏帧序列也能产生高质量分割和平滑网格重建。方法捕获复杂非线性结构的同时保留了显式表示的可解释性和可编辑性，支持组织变形和交互。其设计目的不是单独增加一个模块，而是用上述表示、监督或渲染约束消除原流程在目标任务中的具体失配。提出MedGS框架，基于高斯splatting实现多模态3D医学影像的高保真解剖网格重建。采用多任务设计，高斯原语共享统一几何（位置和时间变形），双色彩分支分别处理帧插值和分割，利用密集图像合成信号正则化几何，在稀疏帧序列下也能实现高质量重建。适用时还应保留论文给出的边界：对切片间间距较大的序列可能需要更多正则化；高斯原语数量对稀疏区域的重建质量有影响。
% METHOD: 2512.11624
\paragraph{Fast and Explicit（医学影像与临床应用）（2025.12）}
Fast and Explicit~\cite{2512_11624} 关注的问题是：把 3DGS 用于医学影像或临床场景 Fast and Explicit 是宽召回补全中发现的论文，主要根据标题、摘要和 arXiv 类别判定其与 3DGS 的关系。
相对原始 3DGS，Fast and Explicit 的变化不是只替换损失函数，而是把表示、初始化、约束或网络预测路径重新组织到任务流程中：该工作被召回是因为论文文本包含 Gaussian/3D Gaussian 相关表述；当前分类先确定其任务域和与 3DGS 的相关性，后续证据流水线再补代码、企业和引用信息。
\textbf{实现细节：} 具体实现可以按输入、核心模块和渲染/优化路径来读。输入端首先取得任务所需的图像、相机、文本、时间序列、语义或几何线索，并把这些线索转换成可以约束高斯中心、尺度、旋转、不透明度、颜色或附加特征的条件。核心结构中，方法使用的主要机制是：从摘要层面看，Fast and Explicit 围绕“把 3DGS 用于医学影像或临床场景”展开。若其正文或代码后续可读，应进一步核对它具体改动的是表示、初始化、渲染、优化、损失、动态建模、压缩或应用接口中的哪一环。 训练或推理时，这些模块通常先给出初始化或中间特征，再通过可微 Gaussian rasterization、重建损失、感知/语义约束、几何正则或任务损失回传到高斯参数；如果方法包含前馈网络，则网络输出的是高斯属性、变形场、压缩码、语义特征或控制变量，而不是只输出最终 2D 图像。渲染路径仍以显式高斯投影、深度排序和 alpha 合成为基础，只是在采样、属性预测、局部编辑、动态变形、压缩编码或下游任务接口处加入额外约束。这样写入本章时，它与前后方法的关系是：相对原始 3DGS 和本章已有锚点方法，它可以理解为 3DGS extends\_or\_uses 宽召回分类阶段的关系判断，后续需要按论文方法部分复核。
\textbf{证据与边界：} 本轮召回分类尚未逐项核对实验表格；后续证据补全阶段需要读取论文源码/PDF或项目页确认数据集、指标和对比方法。 该条来自宽召回批次；源码下载受 arXiv e-print 网络限制，部分论文暂时只能依据标题、摘要和类别分类，需要在入可信书前做人工/LLM精读复核。 本书只采用本地 paper pack、数据库分类、源码/引用登记和论文摘要中可核验的信息；没有本地证据支撑的速度、排行榜或开源完整性不在正文中强化。\textbf{方法信息：} 代码：\url{https://github.com/m-dannecker/gaussian-primitives-for-fast-svr}。
\textbf{设计动机与边界：} 相对原始 3DGS 以独立高斯和图像重建误差为主的设定，Fast and Explicit 把问题落在以下可实现的改动上：从摘要层面看，Fast and Explicit 围绕“把 3DGS 用于医学影像或临床场景”展开。若其正文或代码后续可读，应进一步核对它具体改动的是表示、初始化、渲染、优化、损失、动态建模、压缩或应用接口中的哪一环。其设计目的不是单独增加一个模块，而是用上述表示、监督或渲染约束消除原流程在目标任务中的具体失配。该工作被召回是因为论文文本包含 Gaussian/3D Gaussian 相关表述；当前分类先确定其任务域和与 3DGS 的相关性，后续证据流水线再补代码、企业和引用信息。适用时还应保留论文给出的边界：该条来自宽召回批次；源码下载受 arXiv e-print 网络限制，部分论文暂时只能依据标题、摘要和类别分类，需要在入可信书前做人工/LLM精读复核。
% METHOD: 2601.22990
\paragraph{Self-Supervised Slice-to-Volume Reconstruction with Gaussian Representations for Fetal MRI（医学影像与临床应用）（2026.01）}
Self-Supervised Slice-to-Volume Reconstruction with Gaussian Representations for Fetal MRI~\cite{2601_22990} 关注的问题是：把 3DGS 用于医学影像或临床场景 Self-Supervised Slice-to-Volume Reconstruction with Gaussian Representations for 是宽召回补全中发现的论文，主要根据标题、摘要和 arXiv 类别判定其与 3DGS 的关系。
相对原始 3DGS，Self-Supervised Slice-to-Volume Reconstruction with Gaussian Representations for Fetal MRI 的变化不是只替换损失函数，而是把表示、初始化、约束或网络预测路径重新组织到任务流程中：该工作被召回是因为论文文本包含 Gaussian/3D Gaussian 相关表述；当前分类先确定其任务域和与 3DGS 的相关性，后续证据流水线再补代码、企业和引用信息。
\textbf{实现细节：} 具体实现可以按输入、核心模块和渲染/优化路径来读。输入端首先取得任务所需的图像、相机、文本、时间序列、语义或几何线索，并把这些线索转换成可以约束高斯中心、尺度、旋转、不透明度、颜色或附加特征的条件。核心结构中，方法使用的主要机制是：从摘要层面看，Self-Supervised Slice-to-Volume Reconstruction with Gaussian Representations for 围绕“把 3DGS 用于医学影像或临床场景”展开。若其正文或代码后续可读，应进一步核对它具体改动的是表示、初始化、渲染、优化、损失、动态建模、压缩或应用接口中的哪一环。 训练或推理时，这些模块通常先给出初始化或中间特征，再通过可微 Gaussian rasterization、重建损失、感知/语义约束、几何正则或任务损失回传到高斯参数；如果方法包含前馈网络，则网络输出的是高斯属性、变形场、压缩码、语义特征或控制变量，而不是只输出最终 2D 图像。渲染路径仍以显式高斯投影、深度排序和 alpha 合成为基础，只是在采样、属性预测、局部编辑、动态变形、压缩编码或下游任务接口处加入额外约束。这样写入本章时，它与前后方法的关系是：相对原始 3DGS 和本章已有锚点方法，它可以理解为 3DGS extends\_or\_uses 宽召回分类阶段的关系判断，后续需要按论文方法部分复核。
\textbf{证据与边界：} 本轮召回分类尚未逐项核对实验表格；后续证据补全阶段需要读取论文源码/PDF或项目页确认数据集、指标和对比方法。 该条来自宽召回批次；源码下载受 arXiv e-print 网络限制，部分论文暂时只能依据标题、摘要和类别分类，需要在入可信书前做人工/LLM精读复核。 本书只采用本地 paper pack、数据库分类、源码/引用登记和论文摘要中可核验的信息；没有本地证据支撑的速度、排行榜或开源完整性不在正文中强化。\textbf{方法信息：} 代码：\href{https://github.com/yinsong0510/gaussiansvr-self-supervised-slice-to-volume-reconstruction-with-gaussian-representations}{GaussianSVR 代码仓库}。


医学应用中的 3DGS 价值来自“按模态改造表示”。内窥镜方法围绕动态组织、器械遮挡、高频细节和光照变化改造 4DGS；R$^2$-Gaussian 把高斯变成 X 射线密度核，是 CT/X-ray 路线的核心锚点，DIGS 和 IXGS 分别向动态 CBCT 与真实术中 X-ray 扩展；Deform3DGS 和 G-SHARP 强调手术场景的快速训练、mask 监督和硬件部署；NLI4VolVis 与 ClipGS-VR 把高斯变成医学体数据的交互载体；PINGS-X 则把高斯 splatting 改造成物理场超分辨率近似器。读这一章时，最重要的是不要把所有方法都理解成“医学版 3DGS”，而要看它们各自遵守了哪一种成像物理或临床交互约束。

\textbf{设计动机与边界：} 相对原始 3DGS 以独立高斯和图像重建误差为主的设定，Self-Supervised Slice-to-Volume Reconstruction with Gaussian Representations for Fetal MRI 把问题落在以下可实现的改动上：从摘要层面看，Self-Supervised Slice-to-Volume Reconstruction with Gaussian Representations for 围绕“把 3DGS 用于医学影像或临床场景”展开。若其正文或代码后续可读，应进一步核对它具体改动的是表示、初始化、渲染、优化、损失、动态建模、压缩或应用接口中的哪一环。其设计目的不是单独增加一个模块，而是用上述表示、监督或渲染约束消除原流程在目标任务中的具体失配。该工作被召回是因为论文文本包含 Gaussian/3D Gaussian 相关表述；当前分类先确定其任务域和与 3DGS 的相关性，后续证据流水线再补代码、企业和引用信息。适用时还应保留论文给出的边界：该条来自宽召回批次；源码下载受 arXiv e-print 网络限制，部分论文暂时只能依据标题、摘要和类别分类，需要在入可信书前做人工/LLM精读复核。
\textbf{尾项实现澄清：} GaussianSVR 的动机是胎儿 MRI 切片会受运动破坏，传统 SVR 耗时且需要多组正交切片，监督式方案又依赖不可得的真值体。它以三维高斯表示目标体积，用模拟的前向切片采集模型把渲染切片与实际切片对齐，从而以自监督重建损失训练；多分辨率阶段联合更新高斯参数和空间变换。推理得到体积而非普通新视图，关键约束来自切片再采样一致性，摘要未提供更细的损失权重。
